\section{Conclusions}
\label{sec:conclusions}

This paper revisits a frozen multi-layer assessment of a research-oriented
cross-shard commit simulator in order to distinguish executable observations
from claims that the evaluated models and tests cannot establish. The
original protocol scheduled 1272 tasks: 112 warm-ups and 1160 measured
executions. RQ1 obtained 350 completed bounded property checks without an
observed counterexample, while 70 TLC-large runs were censored by timeout.
The completed base TLC timeout-release checks are vacuous with timeout
disabled.
RQ2 detected ten intentionally introduced \emph{formal-model} mutants; its
perfect score is restricted to selected guard-level defects. RQ3 accepted
300 projected traces originating from deterministic Java-simulator cases
and rejected 300 deliberately corrupted \emph{abstract} traces. The latter
are negative replay controls, not Java implementation failures. RQ4
characterizes observed model-checking costs within finite resource budgets,
not a law of asymptotic growth or tool superiority.

A separate post-campaign Java study (G3-v1) exercised ten code mutations.
Its 110 test executions comprised ten successful controls and 100 mutated
tests (51 FAIL and 49 PASS). All mutated focal tests failed, but their
scenarios, internal checks, and instruments differed. In particular, two
recipient-ownership defects passed every original historical/G1/G2 test in
the battery and became observable after adding a beneficiary-distinguishing
scenario. The existing G2 ledger-journal audit also reported divergence on the same
mutated attempts. This is evidence of insufficient scenario coverage, not proof of
superior detection by a new algorithm.

At the model level, V1 collapses receipt consumption and destination credit,
whereas V2.2 represents more intermediate effects but excludes only selected
contradictory terminal states. G2-P3 documents a Java debit-then-rollback
history that is not represented step-for-step by the examined V2.2 decision
sequence. Accordingly, none of the experiments demonstrates general
Java--V2.2 refinement, durable crash recovery, Byzantine safety or
atomicity across independently failing shard nodes.

The resulting contribution is a reproducible, bounded \emph{negative case
study} of claim adequacy. It shows why an apparently successful collection of
formal checks, mutation outcomes, trace replays and reproduced artifacts can
still lack the observations necessary for specific end-to-end guarantees.
This is not a new trace-validation method; a claim of broad generality would
require an independent implementation, paired scenarios and prospective
comparators. The audit matched the historical runs to their archived records
and regenerated the analysis from those records. The results retain their
original model, workload and observation boundaries.
